% TEMPLATE-MILC-v3.tex - plantilla maestra para todas las guias MILC v3
% Ver scripts/generadores/build-guia-milc.py para la lista completa de
% claves esperadas. El script valida que no queden placeholders sin
% reemplazar antes de escribir el archivo final.

\documentclass[11pt,letterpaper]{article}
\usepackage[margin=1.75cm]{geometry}
\usepackage[table]{xcolor}
\usepackage{fontspec}
\usepackage[spanish]{babel}
\usepackage{tikz}
\usepackage[most]{tcolorbox}
\usepackage{tabularx}
\usepackage{array}
\usepackage{fancyhdr}
\usepackage{titlesec}
\usepackage{ragged2e}
\usepackage{enumitem}
\usepackage{microtype}

\setmainfont{Helvetica Neue}
\setsansfont{Helvetica Neue}
\setlength{\parindent}{0pt}
\setlength{\parskip}{6pt}
\hyphenpenalty=200
\exhyphenpenalty=200
\pretolerance=400
\tolerance=2000
\setlength{\emergencystretch}{1em}
\renewcommand{\arraystretch}{1.25}
\setlist{leftmargin=5mm,itemsep=1mm,topsep=1mm}
\newcolumntype{Y}{>{\RaggedRight\arraybackslash}X}

\definecolor{milcMagenta}{HTML}{E6007E}
\definecolor{milcVino}{HTML}{5A0038}
\definecolor{milcTurquesa}{HTML}{0093A5}
\definecolor{milcVerde}{HTML}{58A923}
\definecolor{milcMostaza}{HTML}{E5B400}
\definecolor{milcOcre}{HTML}{B86600}
\definecolor{milcNegro}{HTML}{241816}
\definecolor{milcGris}{HTML}{F7F7F4}
\definecolor{milcLinea}{HTML}{E8E2DD}
\definecolor{milcGradePrimary}{HTML}{0066FF}
\definecolor{milcGradeDark}{HTML}{003D99}
\definecolor{milcGradeSoft}{HTML}{D6E8FF}
\definecolor{milcGradeAccent}{HTML}{CFE7FF}
\definecolor{milcGradeText}{HTML}{FFFFFF}
\color{milcNegro}

\pagestyle{fancy}
\fancyhf{}
\lhead{\small Tecnología e Informática · Grado 6}
\rhead{\small Guía 19 · MILC v3}
\cfoot{\small\thepage}
\renewcommand{\headrulewidth}{0pt}

\newtcolorbox{titlebox}[1]{
  enhanced,colback=#1,colframe=#1,arc=7mm,boxrule=0pt,
  left=7mm,right=7mm,top=3mm,bottom=3mm,
  before skip=8pt,after skip=8pt,
  fontupper=\sffamily\bfseries\Large\color{white}
}

\newtcolorbox{softbox}[3]{
  enhanced,colback=#2,colframe=#1,borderline west={4pt}{0pt}{#1},
  arc=2mm,boxrule=.4pt,left=4mm,right=4mm,top=3mm,bottom=3mm,
  before skip=6pt,after skip=8pt,title={#3},coltitle=white,
  colbacktitle=#1,fonttitle=\sffamily\bfseries,fontupper=\RaggedRight,
  attach boxed title to top left={xshift=3mm,yshift=-2mm},
  boxed title style={arc=2mm, boxrule=0pt}
}

\newtcolorbox{drawbox}[2]{
  enhanced,colback=white,colframe=#1,arc=4mm,boxrule=1pt,
  left=5mm,right=5mm,top=4mm,bottom=4mm,before skip=8pt,after skip=8pt,
  title={#2},coltitle=white,colbacktitle=#1,fonttitle=\sffamily\bfseries,
  fontupper=\RaggedRight,
  attach boxed title to top left={xshift=4mm,yshift=-2mm},
  boxed title style={arc=3mm, boxrule=0pt}
}

\newtcolorbox{infoband}[1]{
  enhanced,colback=#1!8,colframe=#1!50,arc=2mm,boxrule=.5pt,
  left=4mm,right=4mm,top=2mm,bottom=2mm,before skip=4pt,after skip=4pt,
  fontupper=\small\RaggedRight
}

% Puente narrativo entre secciones (contrato editorial Conéctate).
% Da continuidad: "Acabas de X. Ahora vas a Y porque Z."
% Visualmente: banda izquierda en color del grado, texto en itálica pequeña.
\newtcolorbox{puentebox}{
  enhanced,colback=milcGradeSoft,colframe=milcGradeDark,
  borderline west={3pt}{0pt}{milcGradeDark},
  arc=0mm,boxrule=0pt,left=5mm,right=4mm,top=2.5mm,bottom=2.5mm,
  before skip=4pt,after skip=8pt,
  fontupper=\itshape\small\color{milcNegro}\RaggedRight
}
\newcommand{\puente}[1]{%
  \begin{puentebox}%
    \textcolor{milcGradeDark}{\textbf{\textrightarrow}}\hspace{2mm}#1%
  \end{puentebox}%
}

\newcommand{\linead}[1]{\noindent\color{milcLinea}\rule{#1}{0.5pt}\color{milcNegro}}
\newcommand{\respuesta}[1]{\par\vspace{2mm}\foreach \n in {1,...,#1}{\noindent\color{milcLinea}\rule{\linewidth}{0.5pt}\par\vspace{5mm}}\color{milcNegro}}
\newcommand{\checkbox}{\textcolor{milcGradeDark}{\fbox{\phantom{x}}}\hspace{5pt}}

\newcommand{\phasecard}[4]{
\begin{tcolorbox}[enhanced,colback=#3,colframe=#2,arc=3mm,boxrule=.5pt,height=3.05cm,valign=center,halign=center,left=1.5mm,right=1.5mm,top=2mm,bottom=2mm]
{\sffamily\bfseries\fontsize{22}{24}\selectfont\textcolor{#2}{#1}}\par
{\sffamily\bfseries\small #4}
\end{tcolorbox}
}

\begin{document}
\thispagestyle{empty}
\begin{tikzpicture}[remember picture,overlay]
  \fill[white] (current page.south west) rectangle (current page.north east);
  \path[fill=milcGradePrimary,rounded corners=16mm]
    ([xshift=.55cm,yshift=-.55cm]current page.north west) rectangle
    ([xshift=-.55cm,yshift=3.65cm]current page.south east);

  \node[anchor=north west,text=milcGradeText] at ([xshift=2.0cm,yshift=-1.85cm]current page.north west)
    {\sffamily\bfseries\fontsize{14}{16}\selectfont GRADO};
  \node[anchor=north west,text=milcGradeText,opacity=.82] at ([xshift=2.0cm,yshift=-2.75cm]current page.north west)
    {\sffamily\bfseries\fontsize{9}{11}\selectfont SERIE GUÍAS · TECNOLOGÍA E INFORMÁTICA};

  \begin{scope}[shift={([xshift=-3.05cm,yshift=-2.55cm]current page.north east)}]
    \fill[white,opacity=.80] (-.62,-.52) rectangle (.62,.52);
    \draw[milcGradeText,line width=1pt] (-.62,-.52) rectangle (.62,.52);
    \fill[milcGradeDark] (-.42,-.38) rectangle (-.22,.06);
    \fill[milcGradeAccent] (-.10,-.38) rectangle (.10,.24);
    \fill[milcGradePrimary!60!black] (.22,-.38) rectangle (.42,.38);
  \end{scope}

  \node[anchor=west,text=milcGradeText] at ([xshift=1.9cm,yshift=-12.85cm]current page.north west)
    {\sffamily\bfseries\fontsize{124}{124}\selectfont 6};
  \draw[milcGradeText,line width=1.7pt]
    ([xshift=4.52cm,yshift=-11.68cm]current page.north west) circle (.13cm);

  \node[anchor=west,text=milcGradeText,text width=8.7cm,align=left] at ([xshift=10.35cm,yshift=-11.6cm]current page.north west)
    {
     {\sffamily\bfseries\fontsize{19}{22}\selectfont Postura\\y vida útil ---\\cuidarte tú\\mientras\\usas el equipo}\\[4mm]
     {\sffamily\bfseries\fontsize{23}{26}\selectfont Guía 19-6-TIC}\\[2mm]
     {\sffamily\fontsize{13}{16}\selectfont Período 2 · Hardware y software}\\[5mm]
     {\sffamily\bfseries\fontsize{11}{13}\selectfont Institución Educativa Sor María Juliana}\\[1mm]
     {\sffamily\fontsize{10}{12}\selectfont Autor: Dr. Álvaro Cárdenas Orozco}
    };

  \path[fill=milcGradeSoft,rounded corners=7mm]
    ([xshift=1.35cm,yshift=.85cm]current page.south west) rectangle
    ([xshift=-1.35cm,yshift=4.25cm]current page.south east);
  \draw[milcGradeDark,line width=.65pt,rounded corners=7mm]
    ([xshift=1.35cm,yshift=.85cm]current page.south west) rectangle
    ([xshift=-1.35cm,yshift=4.25cm]current page.south east);
  \node[anchor=center,text=milcNegro,text width=17.0cm,align=center] at ([yshift=2.55cm]current page.south)
    {\sffamily\fontsize{10}{12}\selectfont Metodología MILC · Apertura ancestral · Pensamiento computacional · Triángulo Dussel-Estoicismo-Floridi\\
    \textcolor{milcGradeDark}{\bfseries Producto:} Una \textbf{auditoría de tu propio puesto de trabajo} en el cuaderno: \textbf{(a) dibujo de cómo TE SIENTAS HOY} cuando usas el computador (sé honesto), con anotaciones de qué está bien y qué está mal; \textbf{(b) dibujo de la postura ideal} con etiquetas de los 8 puntos clave (espalda, hombros, codos, muñecas, mirada, distancia, iluminación, pies); \textbf{(c) lista de 5 ajustes concretos} que vas a hacer en tu puesto esta semana. Más \textbf{una promesa firmada} de aplicar la regla \emph{``20-20-20''} todos los días.};
\end{tikzpicture}
\mbox{}
\newpage

\section*{Apertura: Postura y vida útil --- cuidarte tú mientras usas la máquina}

\begin{softbox}{milcGradeDark}{milcGradeSoft}{Saber ancestral}
\textbf{Don Lucho el relojero llegó a los 75 años con vista clara y manos firmes.} ¿Cómo? No fue suerte. Don Lucho \textbf{cuidaba su cuerpo} como cuidaba los relojes. En su vitrina de la calle 14 de Cartago, tenía \textbf{una silla de espaldar alto} (no banco), una \textbf{lámpara de buena luz} encima del mesón, y cada 30 minutos se paraba a \textbf{caminar 2 cuadras} y mirar a lo lejos. Decía: \emph{``Si la vista trabaja todo el día en lo cercano, después no ve a su nieto del otro lado de la plaza.''}. Sus aprendices, en cambio, trabajaban encorvados sobre los relojes 8 horas seguidas, mal sentados, mala luz. Algunos terminaron con problemas de espalda, otros con vista cansada. \textbf{La diferencia entre dominar el oficio toda la vida o gastarse joven está en cómo cuidas tu cuerpo}. Hoy con los computadores pasa lo mismo. Tú tienes 12 años: puedes terminar a los 30 con espalda buena o con espalda dañada, dependiendo de cómo te sientes hoy frente a la pantalla.
\end{softbox}

\begin{softbox}{milcTurquesa}{milcGris}{Saber contemporáneo}
La \textbf{ergonomía} es la ciencia de adaptar la herramienta al cuerpo, no el cuerpo a la herramienta. Cuando usas un computador mal posicionado durante años, el cuerpo se \textbf{daña} en formas que en tu edad no se notan: dolor de cuello, espalda baja, muñecas, ojos cansados, dolor de cabeza por luz mala. \textbf{El daño no llega de un día, llega por sumar horas mal sentadas durante meses y años}. Lo bueno: con 5 reglas simples y la regla \emph{``20-20-20''} (cada 20 minutos miras 20 segundos a 20 pies de distancia), evitas el 80\% del daño. Esta clase es de las más importantes del periodo porque \textbf{las máquinas se reemplazan; tu cuerpo no}.
\end{softbox}

\begin{softbox}{milcMostaza}{milcGris}{Pregunta-puente}
\textit{Si tú a los 35 años tienes la espalda dañada y no puedes trabajar 8 horas, ¿\textbf{de quién es la culpa}? ¿Del computador, del trabajo, o de tu yo de 12 años que se sentaba mal y nadie le enseñó otra cosa?}
\end{softbox}

\begin{softbox}{milcVerde}{milcGris}{Saber y saber hacer}
Al terminar la clase: (1) podrás \textbf{identificar} 8 errores comunes de postura; (2) sabrás \textbf{explicar} las 5 reglas de ergonomía; (3) podrás \textbf{aplicar} la regla 20-20-20; (4) habrás \textbf{evaluado} tu propio puesto y propuesto 5 ajustes concretos.
\end{softbox}

\small
\begin{tabularx}{\textwidth}{>{\bfseries}p{3.0cm}Y}
\rowcolor{milcGradePrimary!12}Autor & Dr. Álvaro Cárdenas Orozco\\
DBA & Identifica componentes y sistemas tecnológicos en su entorno (MEN, T\&I 6°)\\
Referentes & Saber ancestral del relojero · Sistemas como cadena de oficios · Diagnóstico paso a paso\\
Producto final & Una \textbf{auditoría de tu propio puesto de trabajo} en el cuaderno: \textbf{(a) dibujo de cómo TE SIENTAS HOY} cuando usas el computador (sé honesto), con anotaciones de qué está bien y qué está mal; \textbf{(b) dibujo de la postura ideal} con etiquetas de los 8 puntos clave (espalda, hombros, codos, muñecas, mirada, distancia, iluminación, pies); \textbf{(c) lista de 5 ajustes concretos} que vas a hacer en tu puesto esta semana. Más \textbf{una promesa firmada} de aplicar la regla \emph{``20-20-20''} todos los días.\\
\end{tabularx}
\normalsize

\puente{Hoy haces 4 cosas: identificas errores típicos, aprendes las 5 reglas, evaluás tu puesto con honestidad, y firmas la promesa de cambio.}

\begin{titlebox}{milcGradeDark}
Ruta MILC: así vamos a aprender
\end{titlebox}
\begin{tabularx}{\textwidth}{>{\centering\arraybackslash}X>{\centering\arraybackslash}X>{\centering\arraybackslash}X>{\centering\arraybackslash}X}
\phasecard{1}{milcVerde}{milcGris}{Escuta\\[-1mm]\small Observo, escucho y pregunto.} &
\phasecard{2}{milcTurquesa}{milcGris}{Sistematización\\[-1mm]\small Pensamiento computacional.} &
\phasecard{3}{milcMagenta}{milcGris}{Praxis\\[-1mm]\small Construyo y aplico.} &
\phasecard{4}{milcVino}{milcGris}{Evaluación\\[-1mm]\small Reflexión liberadora.}
\end{tabularx}


\puente{Empezamos analizando 8 errores típicos de postura que tal vez tú cometes.}

\begin{titlebox}{milcVerde}
Fase 1 · Escuta
\end{titlebox}

\begin{softbox}{milcVerde}{milcGris}{Escena para escuchar}
\textbf{Actividad 1 · IDENTIFICA --- 8 errores típicos de postura} (10 min · individual). Lee los 8 errores y \textbf{en el cuaderno marca cuáles cometes TÚ regularmente}. Sé honesto, esto es para mejorar. \textbf{Error 1:} encorvar la espalda hacia adelante para acercarte a la pantalla. \textbf{Error 2:} cuello inclinado hacia abajo viendo el celular muchas horas. \textbf{Error 3:} usar el portátil en la cama o sobre las piernas. \textbf{Error 4:} muñecas dobladas hacia arriba o hacia abajo al teclear. \textbf{Error 5:} pantalla muy cerca (a menos de 50 cm). \textbf{Error 6:} luz muy fuerte detrás de la pantalla (reflejo en los ojos). \textbf{Error 7:} sin descansos: 2 horas seguidas sin pararse. \textbf{Error 8:} pies colgando (silla muy alta) o sentado en banco sin espaldar. \emph{Cuenta cuántos cometes}.
\end{softbox}

\checkbox Leí los 8 errores con calma.\\
\checkbox Marqué con honestidad los que cometo yo.\\
\checkbox Conté cuántos errores cometo en total.

\begin{infoband}{milcVerde}
\textbf{Lo que va al cuaderno (Actividad 1 · IDENTIFICA).} Título: «Actividad 1 --- Mis errores de postura». \textbf{Formato:} lista numerada del 1 al 8 con marca SÍ/NO al lado. Al final, escribe el conteo: \emph{``Cometo X errores''}. \textbf{Extensión:} 8 líneas + total. \textbf{Sabes que terminaste cuando} eres honesto contigo mismo en el conteo.
\end{infoband}

\puente{Después aprendes las 5 reglas de ergonomía + la regla 20-20-20.}

\begin{titlebox}{milcTurquesa}
Fase 2 · Sistematización con pensamiento computacional
\end{titlebox}

\begin{softbox}{milcTurquesa}{milcGris}{Cuatro pilares aplicados al tema}
Si cometes \textbf{0-2 errores}, vas muy bien. Si cometes \textbf{3-5}, hay que mejorar. Si cometes \textbf{6-8}, tu cuerpo te va a pasar factura más rápido de lo que crees. Ahora aprendes las \textbf{5 reglas de ergonomía} + la regla \emph{20-20-20} para corregir lo que descubriste.
\end{softbox}

\small
\begin{tabularx}{\textwidth}{>{\bfseries\RaggedRight\arraybackslash}p{3.6cm}Y}
\rowcolor{milcTurquesa!90}\color{white}Pilar & \color{white}\bfseries Aplicación al tema\\
1. Descomposición & \textbf{Regla 1 --- Espalda recta, hombros relajados.} Apoya la espalda en el respaldo de la silla. Hombros caídos (no encogidos hacia las orejas). Si tu silla no tiene respaldo, considera cambiarla --- es la pieza más importante. \textbf{Regla 2 --- Pantalla a la altura de los ojos.} La parte superior de la pantalla debe quedar al mismo nivel que tus ojos cuando miras de frente. Si miras hacia abajo, te dañas el cuello. Si miras hacia arriba, también. Apoya el portátil sobre libros si está muy bajo.\\
2. Reconocimiento de patrones & \textbf{Regla 3 --- Distancia de 50-70 cm a la pantalla.} Aproximadamente un brazo extendido. Más cerca cansas la vista. Más lejos forzas la lectura. \textbf{Regla 4 --- Codos a 90 grados, muñecas rectas.} El antebrazo paralelo al suelo. Las muñecas \emph{nunca} dobladas (ni hacia arriba ni hacia abajo). Si está doblada, ajusta la altura de la silla o pon un cojín. Doblar muñecas durante años causa \emph{síndrome del túnel carpiano} (dolor crónico).\\
3. Abstracción & \textbf{Regla 5 --- Pies en el suelo.} Tus pies deben tocar el piso firmes. Si la silla está muy alta, usa un \emph{apoyo} (un libro grueso, una caja). Pies colgando = mala circulación = piernas pesadas. \textbf{La regla 20-20-20 (¡la más importante!):} cada \emph{20 minutos} de uso, mira durante \emph{20 segundos} a algo que esté a \emph{20 pies} (6 metros) de distancia. Por ejemplo: mira por la ventana al árbol del patio. Esto \textbf{descansa los músculos del ojo} que se fatigan enfocando lo cercano.\\
4. Algoritmo conceptual & \textbf{Iluminación: la regla del dorso.} La luz NO debe estar detrás de la pantalla (te encandila). La luz NO debe estar detrás de ti (refleja en la pantalla). \textbf{La luz debe venir de un lado}, ojalá del techo o de una lámpara lateral. Una buena luz reduce la fatiga visual en un 50\%. \textbf{Descansos activos cada hora:} cada \emph{60 minutos}, párate, camina 2-3 minutos, estira la espalda y los brazos, toma agua. Esto resetea la postura y la circulación.\\
\end{tabularx}
\normalsize

\begin{drawbox}{milcTurquesa}{5 reglas + regla 20-20-20}
\textbf{R1.} Espalda recta, hombros caídos. \\[1mm]
\textbf{R2.} Pantalla a altura de ojos. \\[1mm]
\textbf{R3.} Distancia 50-70 cm a pantalla. \\[1mm]
\textbf{R4.} Codos 90° + muñecas rectas. \\[1mm]
\textbf{R5.} Pies en el suelo. \\[1mm]
\textbf{20-20-20:} cada 20 min, 20 segundos mirando a 20 pies (6 m).
\end{drawbox}

\begin{softbox}{milcMostaza}{milcGris}{Errores comunes y su corrección}
\textbf{Mitos sobre postura que escucharás.} (1) \emph{``Soy joven, eso no me afecta''} --- Falso. El daño se acumula desde los 10 años. Lo que no se nota a los 12 te golpea a los 30. (2) \emph{``Si me siento bien, ya está bien''} --- Falso. El cuerpo se acostumbra a malas posturas y deja de avisar. Por eso hay que medirlo con las 5 reglas. (3) \emph{``El portátil sirve igual en cualquier parte''} --- Falso. El nombre lo dice (\emph{portátil}), pero usarlo siempre en la cama te daña la espalda en años. Para trabajo serio, mesa siempre. (4) \emph{``La regla 20-20-20 es exagerada''} --- Falso. Los oftalmólogos del mundo entero la recomiendan. Lo que parece exagerado hoy te salva la vista mañana.
\end{softbox}

\begin{infoband}{milcTurquesa}
\textbf{Lo que va al cuaderno (Actividad 2 · EXPLICA).} Título: «Actividad 2 --- Las 5 reglas + 20-20-20 con dibujo». \textbf{Formato:} 5 reglas numeradas, una por línea, con un dibujo pequeño al lado de cada una (la espalda con flecha, la pantalla con altura, los codos con ángulo, etc.). Al final, dibuja el reloj de la regla 20-20-20 con sus 3 números. \textbf{Extensión:} 5-6 líneas + dibujos. \textbf{Sabes que terminaste cuando} podrías recitar las 5 reglas y la 20-20-20 sin mirar.
\end{infoband}

\puente{Con las reglas claras, evalúas tu propio puesto y propones ajustes.}

\begin{titlebox}{milcMagenta}
Fase 3 · Praxis con pensamiento computacional
\end{titlebox}

\begin{softbox}{milcMagenta}{milcGris}{Construyendo el producto con los cuatro pilares}
\textbf{Actividad 3 · EVALÚA --- Auditoría de mi puesto + plan de ajustes} (28 min · individual). Vas a hacer 4 cosas en una hoja del cuaderno: \textbf{(a)} dibujar cómo te sientas \emph{hoy} cuando usas el computador en casa (sé honesto); \textbf{(b)} dibujar la postura \emph{ideal} con las 5 reglas; \textbf{(c)} listar 5 ajustes concretos que vas a hacer esta semana; \textbf{(d)} firmar la promesa de aplicar la regla 20-20-20.
\end{softbox}

\small
\begin{tabularx}{\textwidth}{>{\bfseries\RaggedRight\arraybackslash}p{3.6cm}Y}
\rowcolor{milcMagenta!90}\color{white}Pilar & \color{white}\bfseries Aplicación al producto\\
Descomposición & \textbf{Paso 1 --- Título y división.} En la parte superior escribe: \textbf{``Auditoría de mi puesto''} y debajo: \emph{``[Tu nombre], año 2026''}. Divide la hoja en \emph{2 mitades verticales}: izquierda y derecha.\\
Patrones & \textbf{Paso 2 --- Mi postura actual (mitad izquierda).} Encabezado: \textbf{``Cómo me siento HOY''}. Debajo dibuja un \emph{monigote} (figura simple de palitos) frente a un computador, en la postura real que tú usas. Sé honesto: si tienes mala espalda, dibújala curva. Si no apoyas los pies, dibújalos colgando. Si tu pantalla está muy baja, dibújala abajo. Al lado del dibujo, en bullets, anota \textbf{qué tienes mal} (al menos 3 cosas).\\
Abstracción & \textbf{Paso 3 --- Postura ideal (mitad derecha).} Encabezado: \textbf{``Cómo debería sentarme''}. Dibuja otro monigote, esta vez con la postura ideal: espalda recta, pantalla al nivel de ojos, codos 90°, pies en el suelo. Etiqueta los \textbf{8 puntos clave} con números: \emph{(1) espalda, (2) hombros, (3) codos, (4) muñecas, (5) mirada, (6) distancia, (7) iluminación, (8) pies}.\\
Algoritmo de armado & \textbf{Paso 4 --- 5 ajustes concretos (debajo de los dibujos).} Lista numerada de 1 a 5. Cada ajuste debe ser \textbf{concreto y medible}, no \emph{``portarme mejor''}. Ejemplos: \emph{(1) Subir el portátil sobre 2 libros para que quede al nivel de mis ojos; (2) Poner una caja debajo de mis pies para que no cuelguen; (3) Mover la lámpara del lado izquierdo (no detrás de la pantalla); (4) Programar alarma en el celular cada 20 minutos para la regla 20-20-20; (5) Hacer pausa de 3 minutos cada hora para caminar y estirar.}.\\
\end{tabularx}
\normalsize

\begin{drawbox}{milcMagenta}{Antes de cerrar tu auditoría}
\checkbox La mitad izquierda tiene mi postura actual con 3 problemas marcados. \\[1mm]
\checkbox La mitad derecha tiene la postura ideal con los 8 puntos etiquetados. \\[1mm]
\checkbox Lista 5 ajustes concretos y medibles. \\[1mm]
\checkbox La promesa 20-20-20 está subrayada con el reloj dibujado. \\[1mm]
\checkbox Está firmado y fechado. \\[1mm]
\checkbox Voy a aplicar los ajustes durante 1 semana mínimo.
\end{drawbox}

\begin{softbox}{milcMagenta}{milcGris}{Plantilla de trabajo}
\textbf{Estructura.} \textit{Arriba:} título + nombre. \textit{Mitad izquierda:} mi postura actual con 3 problemas. \textit{Mitad derecha:} postura ideal con 8 etiquetas. \textit{Abajo de los dibujos:} 5 ajustes numerados. \textit{Centro inferior:} promesa 20-20-20 subrayada con dibujo de reloj. \textit{Esquina:} firma + fecha.
\end{softbox}

\begin{infoband}{milcMagenta}
\textbf{Lo que va al cuaderno (Actividad 3 · EVALÚA).} Título: «Actividad 3 --- Auditoría de mi puesto + 5 ajustes». \textbf{Formato:} 2 dibujos comparativos + lista de ajustes + promesa firmada. \textbf{Extensión:} 1 hoja entera. \textbf{Sabes que terminaste cuando} mirando la hoja sabes exactamente qué cambiar esta semana y cuándo.
\end{infoband}

\puente{El producto es la auditoría con dibujo + lista de ajustes + promesa firmada.}

\begin{titlebox}{milcMostaza}
Producto de la sesión
\end{titlebox}
\textbf{Auditoría de puesto + 5 ajustes + promesa 20-20-20 · firmada}.
\begin{softbox}{milcMostaza}{milcGris}{Criterios de calidad}
(1) El dibujo de postura actual es honesto (no idealizado). (2) El dibujo de postura ideal tiene los 8 puntos etiquetados. (3) Los 5 ajustes son concretos y medibles. (4) La promesa 20-20-20 está subrayada. (5) Está firmado con nombre y fecha.
\end{softbox}

\puente{Antes de salir, miras tus 5 ajustes: ¿son concretos y medibles?}

\begin{titlebox}{milcVino}
Fase 4 · Evaluación liberadora — cinco dimensiones
\end{titlebox}

\begin{softbox}{milcVino}{milcGris}{Sentido de la evaluación}
Evaluar no es solo revisar una nota. Es reconocer qué aprendí, qué cambió en mí, qué emociones manejé, cómo me relaciono con la ciudad y la comunidad, y qué le dejaré a quien viene detrás.
\end{softbox}

\textbf{Mi reflexión liberadora en cinco dimensiones.} Completa cada frase iniciada:

\small
\begin{tabularx}{\textwidth}{>{\bfseries\RaggedRight\arraybackslash}p{3.4cm}Y>{\RaggedRight\arraybackslash}p{4.6cm}}
\rowcolor{milcVino}\color{white}Dimensión & \color{white}\bfseries Mi respuesta (frame starter) & \color{white}\bfseries Algo que descubrí\\
1. Personal & Lo que más cambió en mí fue \linead{6.5cm} & \linead{4.2cm}\\
2. Emocional & La emoción más fuerte fue \linead{2.5cm} y la regulé \linead{3cm} & \linead{4.2cm}\\
3. Ciudadana & En democracia esto importa porque \linead{6cm} & \linead{4.2cm}\\
4. Local (Cartago) & En mi barrio sería útil para \linead{6.5cm} & \linead{4.2cm}\\
5. Intergeneracional & A mi abuela/abuelo le diría \linead{6.5cm} & \linead{4.2cm}\\
\end{tabularx}
\normalsize

\begin{infoband}{milcVino}
\textbf{Para la próxima sustentación.} Vuelve a esta tabla en seis meses. Verás que las respuestas cambian — no porque lo aprendido cambie, sino porque tú cambias. La evaluación liberadora es una conversación contigo a través del tiempo.
\end{infoband}


\puente{Cerramos con 3 voces que te ayudan a entender por qué cuidar tu cuerpo hoy es regalarle décadas de trabajo a tu yo futuro.}

\begin{titlebox}{milcGradeDark}
Triángulo de pensamiento · cierre filosófico
\end{titlebox}

\begin{softbox}{milcVino}{milcGris}{Dussel · lente del nosotros}
\textit{````Cuidar el cuerpo no es egoísmo: es respetar al que viene después --- el yo de 40 años, los hijos, los compañeros.''''}

Cuando cuidas tu postura hoy, no es \emph{``vanidad''}: es responsabilidad. Tu yo de 40 años va a poder cargar a tus hijos, jugar fútbol con ellos, trabajar 8 horas sin dolor. Tu yo de 60 va a tener vista para leer las cartas de tus nietos. \textbf{Todo eso depende de cómo te sientas frente al computador entre los 12 y los 20 años}. No es exagerar: es ciencia médica.

\textbf{¿Estoy cuidando hoy el cuerpo que mi yo de 40 años necesitará para vivir bien?}
\end{softbox}
\linead{\linewidth}\\[1mm]
\linead{\linewidth}

\begin{softbox}{milcMostaza}{milcGris}{Estoicismo · Marco Aurelio (emperador que cuidó su cuerpo para gobernar 20 años) · cuidado interior}
\textit{````El cuerpo es la única casa que tendrás durante toda tu vida. Cuídala como cuidarías una casa de piedra que no se puede reemplazar.''''}

Marco Aurelio entendió algo: \textbf{tu cuerpo es la única cosa que no se cambia}. El computador lo cambias en 5 años, la silla en 10 años, la casa en 20. \emph{El cuerpo te dura desde el día que naces hasta el día que mueres}. Lo que le hagas, ahí queda. Por eso vale más cuidarlo que cuidar cualquier otra cosa.

\textbf{¿Si solo pudiera cuidar UNA cosa de mi vida con esmero, no debería ser mi cuerpo?}
\end{softbox}
\linead{\linewidth}\\[1mm]
\linead{\linewidth}

\begin{softbox}{milcTurquesa}{milcGris}{Floridi · lente de la infoesfera}
\textit{````En la era digital, la salud corporal se gana o se pierde frente a una pantalla. Por eso es política, no solo medicina.''''}

Tu vida adulta va a tener \textbf{mucho computador y celular}. Trabajos, universidad, comunicación con amigos. Si aprendes ergonomía ahora, vas a estar saludable mientras todo eso pasa. Si no, vas a estar enfermo. \emph{El conocimiento de ergonomía es un derecho que pocos tienen}. Hoy lo tienes tú; pásalo a tu familia.

\textbf{¿Quién en mi familia podría beneficiarse hoy mismo de lo que acabo de aprender sobre postura?}
\end{softbox}
\linead{\linewidth}\\[1mm]
\linead{\linewidth}

\begin{titlebox}{milcVino}
Compromiso final
\end{titlebox}

\small
\begin{tabularx}{\textwidth}{>{\RaggedRight\arraybackslash}p{3.0cm}YYY}
\rowcolor{milcVino}\color{white}\bfseries Criterio & \color{white}\bfseries Lo logré & \color{white}\bfseries En proceso & \color{white}\bfseries Necesito apoyo\\
Comprensión del tema & Explico ideas centrales con mis palabras. & Reconozco ideas pero falta precisión. & Necesito repasar conceptos base.\\
Pensamiento computacional & Apliqué los 4 pilares al diseñar mi producto. & Apliqué algunos pilares. & Necesito releer la fase 2.\\
Aplicación y producto & Mi producto cumple los criterios y se puede revisar. & Producto funcional con ajustes pendientes. & Aún en construcción.\\
Reflexión 5 dimensiones & Respondí cada dimensión con honestidad. & Respondí algunas con detalle. & Mis respuestas son muy generales.\\
Triángulo de pensamiento & Conecté las 3 voces con mi trabajo real. & Conecté al menos una. & No relaciono las voces con mi proceso.\\
\end{tabularx}
\normalsize

\textbf{Hoy entendí que:}\\[1mm]
\linead{\linewidth}\\[1mm]
\linead{\linewidth}

\textbf{Mi compromiso para la próxima sesión es:}\\[1mm]
\linead{\linewidth}\\[1mm]
\linead{\linewidth}

\begin{softbox}{milcMostaza}{milcGris}{Cita de despedida · Marco Aurelio}
\textit{``El obstáculo en el camino se vuelve el camino. Lo que estorba al camino se vuelve camino.''}\\[1mm]
Cada error de hoy, bien observado, se convierte en pieza del camino que pisarás mañana.
\end{softbox}

\small
\begin{tabularx}{\textwidth}{>{\bfseries}p{3.6cm}Y}
\rowcolor{milcGradeDark!12}Nombre completo: & \linead{10cm}\\
Fecha: & \linead{4cm} \quad Curso/grupo: \linead{4cm}\\
Mi firma: & \linead{9cm}\\
\end{tabularx}
\normalsize

\begin{infoband}{milcGradeDark}
\textbf{Para la próxima sesión.} Esta semana aplica los 5 ajustes que escribiste. Pon la alarma del 20-20-20. La próxima clase es la última sesión del periodo: \textbf{Mi computador ideal}. Vas a usar todo lo aprendido (hardware, software, mantenimiento, ergonomía) para diseñar el equipo que TÚ necesitarías comprar y justificar por qué.
\end{infoband}

\end{document}
